% Recovery of the native product and its selector, not incorporation of RH.
\subsection{Multiplicative irreducibility and return periods}
\label{primos-retornos}

The distinction between structure, evaluation, and phase also appears in
the multiplicative sector. APP's local product contains both residue and
quotient; iterating it constructs normal forms and factorizations before
their decimal periods are examined. This development places the expression
\emph{prime mode} in a precise domain and relates irreducibility to the
operations already defined, rather than to an isolated coincidence
of periods. Numeration and the product are constructed below.

Let
\[
 \mathcal W_9=\{\varnothing\}\sqcup
 \bigsqcup_{\ell\geq1}\{1,\ldots,9\}^{\ell}.
\]
Sequences are ordered from the least significant digit. The evaluation
\(\operatorname{ev}_9(u)=\sum_j u_j9^j\), with
\(\operatorname{ev}_9(\varnothing)=0\), is bijective onto the nonnegative
integers: for a positive integer, division of \(n-1\) by \(9\)
uniquely determines the first digit and continuation. This is bijective
nonadic numeration; it retains the boundary representative \(9\).

Addition \(\boxplus_\#\) is obtained by applying APP's additive sheet
to the digits at each position and propagating their quotients. If there
is only one digit, it is the provisional residue; a nonzero incoming
quotient is normalized through a second additive operation. The sum of
quotients is transported to the next position. To multiply a sequence by
a digit \(d\), the multiplicative cell produces
\(u_jd=9q_j+r_j\), and the incoming quotient \(c_j\) is incorporated by
\[
 (w_j,c_{j+1})=
 \begin{cases}
 (r_j,q_j),&c_j=0,\\
 \bigl(R^+(r_j,c_j),q_j+q^+(r_j,c_j)\bigr),&c_j>0.
 \end{cases}
\]
Start with \(c_0=0\) and incorporate the final quotient when nonzero.
During addition, \(c_j\leq2\); during multiplication by a digit,
\(c_j\leq9\). Local operations thus receive arguments within
their domain. Denote this second procedure by \(d\odot_\#u\),
with \(d\odot_\#\varnothing=\varnothing\).

For \(v=(v_0,\ldots,v_{m-1})\), the complete product is constructed
through the recurrence
\[
 a_m=\varnothing,\qquad
 a_j=(9\odot_\#a_{j+1})\boxplus_\#(v_j\odot_\#u),
 \qquad u\boxtimes_\#v=a_0.
\]
The rule composes both local sheets. Its evaluation satisfies
\begin{equation}
 \operatorname{ev}_9(u\boxtimes_\#v)
   =\operatorname{ev}_9(u)\operatorname{ev}_9(v).
 \label{eq:primos-entrelazamiento}
\end{equation}
Indeed, each transition preserves the partial sum plus the pending
quotient weighted by the power of \(9\). The recurrence then gives
\(\operatorname{ev}_9(a_j)=9\operatorname{ev}_9(a_{j+1})+
v_j\operatorname{ev}_9(u)\), whose iteration proves
\eqref{eq:primos-entrelazamiento}. Uniqueness of normal form transports
associativity and commutativity to the constructed product; its identity is \((1)\).

\begin{definition}[Reduced multiplicative coproduct]
On the finitely supported vector space generated by nonempty normal
forms, define
\[
 \widetilde\Delta_\# e_w
 =\sum_{\substack{u\boxtimes_\#v=w\\u,v\ne(1)}}e_u\otimes e_v.
\]
A non-unit normal form is irreducible when its image under
\(\widetilde\Delta_\#\) is zero.
\end{definition}

Factorization fibers are finite by
\eqref{eq:primos-entrelazamiento}. In the completion
\(\mathcal H_\#=\ell^2(\mathcal W_9\setminus\{\varnothing\})\),
images of distinct forms have disjoint supports. If \(d_\#(w)\)
counts ordered non-unit factorizations, the domain of the closure is
\[
 \operatorname{Dom}(\overline{\widetilde\Delta_\#})
 =\left\{x\in\mathcal H_\#:
       \sum_w d_\#(w)|x_w|^2<\infty\right\}.
\]
Indeed, the squared norm of the image is this weighted sum. Hence
the operator
\[
 A_\#=(\overline{\widetilde\Delta_\#})^*
                  \overline{\widetilde\Delta_\#}
\]
is diagonal, positive, and self-adjoint: its eigenvalue at \(e_w\)
is the number of ordered non-unit factorizations of \(w\). The projection
\[
 P_{\rm irr}=(I-|e_{(1)}\rangle\langle e_{(1)}|)
                    \mathbf1_{\{0\}}(A_\#)
\]
selects exactly the irreducible forms. Through multiplicative evaluation,
these correspond to ordinary primes. Selection depends on factorization
structure and keeps the unit state separate.

For a value \(n\geq2\) coprime to \(10\), decimal division induces
the residue permutation \(r\mapsto10r\pmod n\). Return of residue
\(1\) has period
\[
 \tau_n=\operatorname{ord}_n(10).
\]
If \(n=\prod_p p^{a_p}\), the Chinese remainder theorem gives
\[
 \tau_n=\operatorname{lcm}_{p^{a_p}\parallel n}
                    \operatorname{ord}_{p^{a_p}}(10).
\]
Composite return synchronizes the returns of prime powers.
For a prime \(p\notin\{2,3,5\}\),
\[
 \tau_p\mid p-1,\qquad
 M_p=\frac{10^{\tau_p}-1}{p}\in\mathbb N,
 \qquad M_p\equiv0\pmod9.
\]
The last congruence expresses nonadic closure of decimal repetition:
\(10^{\tau_p}-1\) is a multiple of \(9\), and \(p\) is invertible
modulo \(9\). The same congruence holds for every \(n\) coprime
to \(90\). In particular, \(7,13,77,91,143\) have decimal period
\(6\), although only the first two are irreducible. The coproduct and
the period therefore record different properties of the same constructed value.

This differentiation specifies the relation with aperiodic refinement.
The vacancy sequence determines increments of the capacity index;
factorization determines multiplicative modes; decimal division
determines their return times. The nonadic representative preserves
phase closure, while quotient and factorization retain information
that this closure does not distinguish. The connection to the treatise's
spectral constructions resides in these operators and their intertwinings.
Proof of their subsequent analytic results retains its own development;
it is not replaced by periodicity of a residue observation.

\subsubsection{Vacancy discrepancy and Dirichlet readings}

The connection also admits an exact analytic expression. Retain the
same coding \(z_n=z_n(0)\), for \(n\geq1\), and consider two
subsequent observables: the series over all indices and the product
over indices selected by irreducibility. For \(\operatorname{Re}s>1\),
\[
 Z_{\rm vac}^{+}(s)=\sum_{n\geq1}\frac{z_n}{n^s},\qquad
 Z_{\rm vac}^{\times}(s)=\prod_p(1-p^{-s})^{-z_p}.
\]
Both converge absolutely in this half-plane. Their common factor is
density \(\delta\), but their discrepancies differ.

\Needspace{21\baselineskip}
\begin{proposition}[Separation of vacancy density]
We have
\[
 Z_{\rm vac}^{+}(s)=\delta\zeta(s)+H_{\rm vac}(s),
 \qquad H_{\rm vac}\in\mathcal O(\operatorname{Re}s>0).
\]
In \(\operatorname{Re}s>1\), with logarithms defined by absolutely
convergent Euler series,
\[
 Z_{\rm vac}^{\times}(s)=\zeta(s)^\delta G_{\rm vac}(s),\qquad
 \log G_{\rm vac}(s)=
 \sum_p(z_p-\delta)\log(1-p^{-s})^{-1}.
\]
\end{proposition}
\begin{proof}
The centered partial sums satisfy
\[
 A_N:=\sum_{n=1}^N(z_n-\delta)
 =\lfloor(N+1)\delta\rfloor-N\delta,\qquad |A_N|<1.
\]
Abel summation represents the centered term by
\(s\int_1^\infty A_{\lfloor x\rfloor}x^{-s-1}\,dx\).
The integral converges locally uniformly in \(\operatorname{Re}s>0\),
proving holomorphy. The second identity follows by separating
\(z_p=\delta+(z_p-\delta)\) in the logarithm of the product,
within its domain of absolute convergence.
\end{proof}

The decomposition reveals what the reading on primes adds: it examines
the same vacancy sequence on a subset determined by factorization.
Controlling discrepancy over all integers gives the first holomorphic
continuation; the second observable preserves its own discrepancy on
primes. These are two applications of the same discrete structure,
with explicit analytic domains. This articulation explains their
relevance to the arithmetic kernel without transferring the proof
of the treatise's terminal results to this article.
